% Tables: what every table holds, decided 17 Nov 2026. Values are "xx" until the runs exist;
% "xx $\pm$ xx" is the mean and standard deviation across the three seeds, "xx (xx to xx)" an
% estimate with its 95% confidence interval from a bootstrap over cases (paired by case for
% differences). Facts that are already fixed (reward weights, dataset counts, the thinking test) are
% filled in.
% Placement (paper/decisions.md, 17 Nov): the main text keeps its six figures; the main results go in
% Extended Data Table 1, which takes one of the ten Extended Data slots (one is still free); the thinking test moves to the
% Supplementary; every other table is a Supplementary Table in the separate PDF.

\newcommand{\xx}{xx}
\newcommand{\msd}{xx $\pm$ xx}
\newcommand{\ci}{xx (xx to xx)}
\newcolumntype{C}{>{\centering\arraybackslash}X}
\newcommand{\grouprow}[2]{\multicolumn{#1}{@{}l}{\textit{#2}}\\}

\section*{Tables (draft, values xx)}

Extended Data Table~1 holds the main results; Supplementary Tables~1--13 hold the rest. Every value
printed as xx is still to be measured.

% ================================================================================================
\renewcommand{\tablename}{Extended Data Table}
\setcounter{table}{0}

\begin{table}[p]
\centering\footnotesize\setlength{\tabcolsep}{3pt}
\caption{\textbf{Main results on the two test sets.} Mean $\pm$ s.d. across three training seeds for
the untrained Qwen3.5-9B and the two models trained from it, with every episode scored on both rewards.
The NeuroBench test holds the 119 synthetic cases of the four conditions kept out of training;
MIMIC-IV holds 100 recorded admissions of nine conditions, 88 of them from conditions seen in training
or validation (Supplementary Table~1 splits them). The difference is policy reward minus outcome
reward, with a 95\% confidence interval from a bootstrap over cases paired by case. $P$ comes from a
two-sided Wilcoxon signed-rank test (McNemar for proportions), adjusted by Benjamini--Hochberg across
the rows of each test set.}
\label{tab:main}
\begin{tabularx}{\linewidth}{@{}lCCCCc@{}}
\toprule
Measure & Untrained & Outcome reward & Policy reward & Difference (95\% CI) & $P$ \\
\midrule
\grouprow{6}{NeuroBench test (119 cases)}
Correct diagnoses (\%) & \msd & \msd & \msd & \ci & \xx \\
Required tests ordered (\%) & \msd & \msd & \msd & \ci & \xx \\
Recommended tests ordered (\%) & \msd & \msd & \msd & \ci & \xx \\
Cases with an unsafe action (\%) & \msd & \msd & \msd & \ci & \xx \\
Useless tests per case & \msd & \msd & \msd & \ci & \xx \\
Tests per case & \msd & \msd & \msd & \ci & \xx \\
Cost per case (US\$) & \msd & \msd & \msd & \ci & \xx \\
Outcome score & \msd & \msd & \msd & \ci & \xx \\
Policy score & \msd & \msd & \msd & \ci & \xx \\
\midrule
\grouprow{6}{MIMIC-IV neurology admissions (100 cases)}
Correct diagnoses (\%) & \msd & \msd & \msd & \ci & \xx \\
Required tests ordered (\%) & \msd & \msd & \msd & \ci & \xx \\
Recommended tests ordered (\%) & \msd & \msd & \msd & \ci & \xx \\
Cases with an unsafe action (\%) & \msd & \msd & \msd & \ci & \xx \\
Useless tests per case & \msd & \msd & \msd & \ci & \xx \\
Tests per case & \msd & \msd & \msd & \ci & \xx \\
Cost per case (US\$) & \msd & \msd & \msd & \ci & \xx \\
Outcome score & \msd & \msd & \msd & \ci & \xx \\
Policy score & \msd & \msd & \msd & \ci & \xx \\
\bottomrule
\end{tabularx}
\end{table}

% ================================================================================================
\renewcommand{\tablename}{Supplementary Table}
\setcounter{table}{0}

\begin{table}[p]
\centering\footnotesize\setlength{\tabcolsep}{3pt}
\caption{\textbf{MIMIC-IV results split by how each condition was used}, the numbers behind Fig.~5b. Twelve admissions belong to
conditions kept out of training (hepatic encephalopathy 10, bacterial meningitis 2). The other 88
belong to conditions seen in training (ischaemic stroke 38, subarachnoid haemorrhage 12, status
epilepticus 9, Guillain--Barr\'e syndrome 3, functional neurological disorder 2, focal epilepsy 2) or
used to choose the checkpoint (multiple sclerosis 22). Mean $\pm$ s.d. across three seeds.}
\label{tab:mimic-split}
\begin{tabularx}{\linewidth}{@{}lCCC@{}}
\toprule
Measure & Untrained & Outcome reward & Policy reward \\
\midrule
\grouprow{4}{Conditions kept out of training (12 cases)}
Correct diagnoses (\%) & \msd & \msd & \msd \\
Required tests ordered (\%) & \msd & \msd & \msd \\
Cases with an unsafe action (\%) & \msd & \msd & \msd \\
Policy score & \msd & \msd & \msd \\
\midrule
\grouprow{4}{Conditions seen in training or validation (88 cases)}
Correct diagnoses (\%) & \msd & \msd & \msd \\
Required tests ordered (\%) & \msd & \msd & \msd \\
Cases with an unsafe action (\%) & \msd & \msd & \msd \\
Policy score & \msd & \msd & \msd \\
\bottomrule
\end{tabularx}
\end{table}

% ------------------------------------------------------------------------------------------------
\begin{table}[p]
\centering\footnotesize\setlength{\tabcolsep}{3pt}
\caption{\textbf{Untrained models on the NeuroBench test}, the numbers behind Fig.~2. Mean $\pm$
s.d. across three seeds. The last column gives each model's rank by correct diagnoses and by policy
score.}
\label{tab:panel}
\begin{tabularx}{\linewidth}{@{}lCCCCCCc@{}}
\toprule
Model & Correct diagnoses (\%) & Required tests ordered (\%) & Cases with an unsafe action (\%) & Useless tests per case & Cost per case (US\$) & Policy score & Rank \\
\midrule
Qwen3.5-9B & \msd & \msd & \msd & \msd & \msd & \msd & \xx/\xx \\
Granite 4.1 8B & \msd & \msd & \msd & \msd & \msd & \msd & \xx/\xx \\
Gemma 4 E4B & \msd & \msd & \msd & \msd & \msd & \msd & \xx/\xx \\
Larger model 1 & \msd & \msd & \msd & \msd & \msd & \msd & \xx/\xx \\
Larger model 2 & \msd & \msd & \msd & \msd & \msd & \msd & \xx/\xx \\
Larger model 3 & \msd & \msd & \msd & \msd & \msd & \msd & \xx/\xx \\
\bottomrule
\end{tabularx}
\end{table}

% ------------------------------------------------------------------------------------------------
\begin{table}[p]
\centering\footnotesize\setlength{\tabcolsep}{3pt}
\caption{\textbf{Results for each condition of the NeuroBench test}, the numbers behind Fig.~3c.
Mean $\pm$ s.d. across three seeds.}
\label{tab:per-condition}
\begin{tabularx}{\linewidth}{@{}lCCC@{}}
\toprule
Measure & Untrained & Outcome reward & Policy reward \\
\midrule
\grouprow{4}{Bacterial meningitis (\xx{} cases)}
Correct diagnoses (\%) & \msd & \msd & \msd \\
Required tests ordered (\%) & \msd & \msd & \msd \\
\grouprow{4}{High-grade glioma (\xx{} cases)}
Correct diagnoses (\%) & \msd & \msd & \msd \\
Required tests ordered (\%) & \msd & \msd & \msd \\
\grouprow{4}{Hepatic encephalopathy (\xx{} cases)}
Correct diagnoses (\%) & \msd & \msd & \msd \\
Required tests ordered (\%) & \msd & \msd & \msd \\
\grouprow{4}{Myasthenia gravis (\xx{} cases)}
Correct diagnoses (\%) & \msd & \msd & \msd \\
Required tests ordered (\%) & \msd & \msd & \msd \\
\bottomrule
\end{tabularx}
\end{table}

% ------------------------------------------------------------------------------------------------
\begin{table}[p]
\centering\footnotesize\setlength{\tabcolsep}{3pt}
\caption{\textbf{Controls and ablations on the NeuroBench test.} The scripted agents mark the ends of
each scale. The agents without reinforcement learning show how much of the gain the training itself
brings. Mean $\pm$ s.d. across three seeds for the sampled agents; the scripted agents are
deterministic.}
\label{tab:controls}
\begin{tabularx}{\linewidth}{@{}>{\raggedright\arraybackslash}XCCCC@{}}
\toprule
Agent & Correct diagnoses (\%) & Required tests ordered (\%) & Outcome score & Policy score \\
\midrule
\grouprow{5}{Scripted}
Perfect: the required work-up, then the policy's diagnosis & \xx & \xx & \xx & \xx \\
Orders every test, then the policy's diagnosis & \xx & \xx & \xx & \xx \\
\grouprow{5}{Untrained Qwen3.5-9B}
No test, diagnosis from the history alone & \msd & \msd & \msd & \msd \\
No test menu, tests requested in free text & \msd & \msd & \msd & \msd \\
Work-up rules written in the prompt & \msd & \msd & \msd & \msd \\
\grouprow{5}{Trained Qwen3.5-9B}
Supervised fine-tuning on approved episodes & \msd & \msd & \msd & \msd \\
Outcome reward without the cost term & \msd & \msd & \msd & \msd \\
Policy reward without the useless-test penalty & \msd & \msd & \msd & \msd \\
\bottomrule
\end{tabularx}
\end{table}

% ------------------------------------------------------------------------------------------------
\begin{table}[p]
\centering\footnotesize\setlength{\tabcolsep}{4pt}
\caption{\textbf{The two rewards.} The policy score (clinical-policy-v3) is a weighted sum of
components between 0 and 1. The process components count in proportion to the required coverage,
and cost and token efficiency in proportion to diagnosis $\times$ required coverage. The outcome reward
trains the second model; $B$ is the 95th percentile of the untrained model's cost on the 419 training
cases.}
\label{tab:rewards}
\begin{tabularx}{\linewidth}{@{}l>{\raggedright\arraybackslash}Xr@{}}
\toprule
Component & Definition & Weight \\
\midrule
\grouprow{3}{Policy reward}
Diagnosis & an accepted label of the policy's diagnosis in the asserted text & 0.26 \\
Required coverage & share of required criteria met, with partial credit per required argument & 0.19 \\
Recommended coverage & share of recommended criteria met & 0.05 \\
Tool accuracy & each criterion matched by its first action only; $\times$ required coverage & 0.11 \\
Safety & no action the policy says to avoid; $\times$ required coverage & 0.13 \\
Waste avoidance & no low-value action; $\times$ required coverage & 0.07 \\
Sequence & ordering constraints between the actions taken; $\times$ required coverage & 0.06 \\
Stopping & no unjustified action after the required set & 0.06 \\
Cost efficiency & reference cost of the required work-up over the actual cost; gated & 0.04 \\
Token efficiency & $1/(1+\text{completion tokens}/2048)$; gated & 0.03 \\
Assessment, ICD-10 code & reported, not weighted & 0 \\
Useless-test penalty & per action that meets no criterion & $-0.03$ \\
Invalid actions & score divided by $1 + 0.1\,n_\text{invalid}$ & -- \\
Caps & contraindicated action 0.10, harmful action 0.35, hard sequence violation 0.50 & -- \\
\midrule
\grouprow{3}{Outcome reward}
Diagnosis & as above & 0.85 \\
Cost & $1 - \text{cost}/B$, with $B = \text{US\$}\,5{,}913$ & 0.15 \\
\bottomrule
\end{tabularx}
\end{table}

% ------------------------------------------------------------------------------------------------
\begin{table}[p]
\centering\footnotesize\setlength{\tabcolsep}{3pt}
\caption{\textbf{Datasets.} NeuroBench holds 600 synthetic cases of 20 conditions, split by whole
condition; two cases are left out because their answer belongs to another split. MIMIC-IV holds 100
recorded emergency admissions, used for testing only; their policies are carried over from the
NeuroBench cases of the same condition.}
\label{tab:data}
\begin{tabularx}{\linewidth}{@{}lCCCCl@{}}
\toprule
Set & Conditions & Cases & Simple / moderate / complex & Policies approved & Use \\
\midrule
NeuroBench training & 14 & 419 & \xx{} / \xx{} / \xx & \xx & training \\
NeuroBench validation & 2 & 60 & \xx{} / \xx{} / \xx & \xx & checkpoint choice \\
NeuroBench test & 4 & 119 & \xx{} / \xx{} / \xx & \xx & test \\
MIMIC-IV neurology & 9 & 100 & -- & \xx & external test \\
\bottomrule
\end{tabularx}
\end{table}

% ------------------------------------------------------------------------------------------------
\begin{table}[p]
\centering\footnotesize\setlength{\tabcolsep}{3pt}
\caption{\textbf{Models and inference settings.} Sampling and the thinking switch follow each model
card and are fixed per model; runs are replicated by seed. Model weights accessed on \xx.}
\label{tab:models}
\begin{tabularx}{\linewidth}{@{}lCCCCCC@{}}
\toprule
Model & Revision & Licence & Temperature & Top-$p$ & Max tokens per turn & Thinking \\
\midrule
Qwen3.5-9B & \xx & \xx & \xx & \xx & \xx & off \\
Granite 4.1 8B & \xx & \xx & \xx & \xx & \xx & \xx \\
Gemma 4 E4B & \xx & \xx & \xx & \xx & \xx & off \\
Larger model 1 & \xx & \xx & \xx & \xx & \xx & \xx \\
Larger model 2 & \xx & \xx & \xx & \xx & \xx & \xx \\
Larger model 3 & \xx & \xx & \xx & \xx & \xx & \xx \\
\bottomrule
\end{tabularx}
\end{table}

% ------------------------------------------------------------------------------------------------
\begin{table}[p]
\centering\footnotesize\setlength{\tabcolsep}{4pt}
\caption{\textbf{Training and compute.} Both models start from the same untrained Qwen3.5-9B and
differ only in the reward. The checkpoint is chosen on the validation cases by each model's own
reward.}
\label{tab:training}
\begin{tabularx}{\linewidth}{@{}>{\raggedright\arraybackslash}XCC@{}}
\toprule
Setting & Outcome reward & Policy reward \\
\midrule
Algorithm & GRPO, DAPO loss & GRPO, DAPO loss \\
Training steps & 500 & \xx \\
Learning rate & $2\times10^{-5}$ & $2\times10^{-5}$ \\
Rollouts per case; cases per step & \xx; \xx & \xx; \xx \\
LoRA rank; $\alpha$ & \xx; \xx & \xx; \xx \\
Checkpoint chosen (step), seeds 1/2/3 & \xx/\xx/\xx & \xx/\xx/\xx \\
GPU; hours per seed & \xx; \xx & \xx; \xx \\
\bottomrule
\end{tabularx}
\end{table}

% ------------------------------------------------------------------------------------------------
\begin{table}[p]
\centering\footnotesize\setlength{\tabcolsep}{4pt}
\caption{\textbf{Agreement between physicians.} Two neurologists reviewed each policy on their own.
Cohen's $\kappa$ per policy dimension on the \xx{} cases reviewed by both, and the agreement of the
blind episode ratings with the policy score.}
\label{tab:agreement}
\begin{tabularx}{\linewidth}{@{}>{\raggedright\arraybackslash}XCC@{}}
\toprule
Dimension & Cases & $\kappa$ (95\% CI) \\
\midrule
Diagnosis & \xx & \ci \\
Required tests & \xx & \ci \\
Recommended tests & \xx & \ci \\
Actions to avoid & \xx & \ci \\
Sequence & \xx & \ci \\
Blind episode rating against policy score (Spearman $\rho$) & \xx & \ci \\
\bottomrule
\end{tabularx}
\end{table}

% ------------------------------------------------------------------------------------------------
\begin{table}[p]
\centering\footnotesize\setlength{\tabcolsep}{4pt}
\caption{\textbf{Cost of running an episode.} Mean $\pm$ s.d. per episode on the NeuroBench test, on
one \xx{} GPU with vLLM \xx.}
\label{tab:runtime}
\begin{tabularx}{\linewidth}{@{}lCCCC@{}}
\toprule
Model & Turns & Completion tokens & Wall time (s) & Test cost (US\$) \\
\midrule
Untrained & \msd & \msd & \msd & \msd \\
Outcome reward & \msd & \msd & \msd & \msd \\
Policy reward & \msd & \msd & \msd & \msd \\
\bottomrule
\end{tabularx}
\end{table}

% ------------------------------------------------------------------------------------------------
% Source: paper/decisions.md (17 Nov) and EOS runs/thinking_val_20261117/. To check before
% submission: whether invalid actions and failed episodes are totals over the 180 episodes or means
% per seed, the exact budgets in tokens, and the s.d. across seeds.
\begin{table}[p]
\centering\footnotesize\setlength{\tabcolsep}{4pt}
\caption{\textbf{Thinking off and on.} Untrained Qwen3.5-9B on the 60 validation cases, three seeds
(180 episodes per row), with thinking off or with a thinking budget per turn. Thinking used about
600 tokens per turn and never more than 6,000, so the budget never bound. With thinking on, the
model often asked for several tests in one turn or answered without ordering any, so the trained
models run with thinking off.}
\label{tab:thinking}
\begin{tabularx}{\linewidth}{@{}lCCCCC@{}}
\toprule
Thinking & Policy score & Correct diagnoses (\%) & Tests per case & Invalid actions & Failed episodes \\
\midrule
Off & 0.41 & 88 & 3.6 & 29 & 1 \\
On, 8k tokens & 0.38 & 82 & 2.9 & 188 & 18 \\
On, 16k tokens & 0.39 & 84 & 2.7 & 183 & 16 \\
On, 32k tokens & 0.38 & 81 & 2.8 & 173 & 12 \\
\bottomrule
\end{tabularx}
\end{table}

% ------------------------------------------------------------------------------------------------
% Added 18 Nov 2026 with Fig. 4c: the comparison with the neurologists left Fig. 2 and its numbers left
% Supplementary Table 2.
\begin{table}[p]
\centering\footnotesize\setlength{\tabcolsep}{3pt}
\caption{\textbf{Models against neurologists who solved the same cases}, the numbers behind Fig.~4c.
\xx{} neurologists solved \xx{} test cases through the same interface as the models. The first two
columns read each model against the neurologist of the case; the last two apply to the neurologists
too. Mean (95\% confidence interval from a bootstrap over cases); the models' values are averaged over
three seeds.}
\label{tab:neurologists}
\begin{tabularx}{\linewidth}{@{}lCCCC@{}}
\toprule
 & Agrees on the diagnosis (\%) & Neurologist's tests ordered (\%) & Required tests ordered (\%) & Tests per case \\
\midrule
Neurologists & -- & -- & \ci & \ci \\
Untrained & \ci & \ci & \ci & \ci \\
Outcome reward & \ci & \ci & \ci & \ci \\
Policy reward & \ci & \ci & \ci & \ci \\
\bottomrule
\end{tabularx}
\end{table}

% ------------------------------------------------------------------------------------------------
% Added 18 Nov 2026 with Fig. 5c.
\begin{table}[p]
\centering\footnotesize\setlength{\tabcolsep}{3pt}
\caption{\textbf{Models against the clinicians who treated the same 100 patients}, the numbers behind
Fig.~5c. The first two columns read each model against the tests recorded in the MIMIC-IV admission;
tests per case applies to the clinicians too and counts, for them, the catalogue tests recorded during
the stay. Mean (95\% confidence interval from a bootstrap over cases); the models' values are averaged
over three seeds.}
\label{tab:clinicians}
\begin{tabularx}{\linewidth}{@{}lCCC@{}}
\toprule
 & Tests the clinicians also did (\%) & Clinicians' tests ordered (\%) & Tests per case \\
\midrule
Clinicians & -- & -- & \ci \\
Untrained & \ci & \ci & \ci \\
Outcome reward & \ci & \ci & \ci \\
Policy reward & \ci & \ci & \ci \\
\bottomrule
\end{tabularx}
\end{table}

